\section{Множества и отношения}
\label{sec:sets-relations}

\subsection{Множества и операции}

Запись $x\in A$ означает принадлежность элемента множеству. Подмножество:
\[
 A\subseteq B
 \quad\Longleftrightarrow\quad
 \forall x\;(x\in A\Rightarrow x\in B).
\]
Два множества равны тогда и только тогда, когда каждое является подмножеством
другого.

Основные операции:
\[
\begin{aligned}
 A\cup B&=\{x:x\in A\text{ или }x\in B\},\\
 A\cap B&=\{x:x\in A\text{ и }x\in B\},\\
 A\setminus B&=\{x:x\in A,\ x\notin B\},\\
 A\triangle B&=(A\setminus B)\cup(B\setminus A).
\end{aligned}
\]
Относительно универсального множества $U$ вводится дополнение
$\overline A=U\setminus A$. Законы де Моргана:
\[
 \overline{A\cup B}=\overline A\cap\overline B,
 \qquad
 \overline{A\cap B}=\overline A\cup\overline B.
\]

Множество всех подмножеств обозначается $\mathcal P(A)$ или $2^A$. Если
$|A|=n$, то
\[
    |\mathcal P(A)|=2^n,
\]
поскольку каждое подмножество задаётся бинарным вектором принадлежности длины
$n$.

\subsection{Декартово произведение и мощности}

Декартово произведение
\[
    A\times B=\{(a,b):a\in A,\ b\in B\}.
\]
Для конечных множеств
\[
    |A\times B|=|A|\,|B|.
\]

Множества $A$ и $B$ равномощны, если между ними существует биекция. Множество
счётно бесконечно, если равномощно $\mathbb N$; ``не более чем счётное'' означает
конечное или счётно бесконечное.

Классические факты:
\begin{itemize}
 \item $\mathbb Z$ и $\mathbb Q$ счётно бесконечны;
 \item конечное декартово произведение счётных множеств счётно;
 \item счётное объединение счётных множеств не более чем счётно;
 \item $\mathbb R$ несчётно.
\end{itemize}
Несчётность $[0,1]$ можно показать диагональным аргументом Кантора: из любого
предполагаемого списка десятичных записей строится число, отличающееся от
$n$-го элемента списка в $n$-й цифре.

\textbf{Теорема Кантора.} Для любого множества $A$ не существует сюръекции
$A\to\mathcal P(A)$; в частности,
\[
    |A|<|\mathcal P(A)|.
\]
Доказательство использует диагональное множество
\[
    D=\{a\in A:a\notin f(a)\}.
\]
Если бы $D=f(d)$, то $d\in D\iff d\notin D$, противоречие.

\subsection{Бинарные отношения}

Бинарное отношение между $A$ и $B$ --- подмножество
\[
    R\subseteq A\times B.
\]
Запись $aRb$ означает $(a,b)\in R$. Область определения и множество значений:
\[
 \operatorname{Dom}R=\{a:\exists b\ aRb\},
 \qquad
 \operatorname{Im}R=\{b:\exists a\ aRb\}.
\]

Обратное отношение:
\[
    R^{-1}=\{(b,a):(a,b)\in R\}.
\]
Композиция $R\subseteq A\times B$, $S\subseteq B\times C$:
\[
 S\circ R=\{(a,c):\exists b\;(aRb\land bSc)\}.
\]
Композиция ассоциативна:
\[
    T\circ(S\circ R)=(T\circ S)\circ R,
\]
но обычно некоммутативна.

\subsection{Свойства отношений и замыкания}

Для $R\subseteq A\times A$:
\begin{itemize}
 \item рефлексивность: $aRa$ для всех $a$;
 \item антирефлексивность: $aRa$ не выполняется ни для одного $a$;
 \item симметричность: $aRb\Rightarrow bRa$;
 \item антисимметричность: $aRb$ и $bRa$ влекут $a=b$;
 \item транзитивность: $aRb$ и $bRc$ влекут $aRc$;
 \item связность (полнота): для различных $a,b$ выполнено $aRb$ или $bRa$.
\end{itemize}

Транзитивное замыкание
\[
    R^+=\bigcup_{k\ge1}R^k
\]
--- наименьшее транзитивное отношение, содержащее $R$. Рефлексивно-транзитивное
замыкание
\[
    R^*=I\cup R^+,
\]
где $I=\{(a,a):a\in A\}$. Для конечного ориентированного графа отношение
$R^*$ совпадает с достижимостью по путям.

\subsection{Отношение эквивалентности и разбиение}

Отношение эквивалентности рефлексивно, симметрично и транзитивно. Класс элемента
$a$:
\[
    [a]=\{x\in A:xRa\}.
\]

\textbf{Теорема.} Классы эквивалентности образуют разбиение $A$: каждый элемент
лежит в некотором классе, а два класса либо совпадают, либо не пересекаются.
Обратно, любое разбиение $A$ на непустые попарно непересекающиеся блоки задаёт
отношение эквивалентности --- два элемента эквивалентны, если принадлежат одному
блоку.

Схема доказательства непересекаемости: если $[a]\cap[b]\ne\varnothing$, выберем
$x$ в пересечении. Из $aRx$ и $xRb$ (с учётом симметрии) следует $aRb$; затем
транзитивность показывает $[a]=[b]$.

Множество классов
\[
    A/R
\]
называется фактор-множеством.

\subsection{Отношения порядка}

Частичный порядок --- рефлексивное, антисимметричное и транзитивное отношение.
Если любые два элемента сравнимы, порядок называется линейным.

Связанный с ним строгий порядок задаётся
\[
    a\prec b\iff a\preceq b\text{ и }a\ne b.
\]
Он антирефлексивен и транзитивен.

Нужно различать:
\begin{itemize}
 \item минимальный и наименьший элементы;
 \item максимальный и наибольший элементы;
 \item нижние/верхние границы;
 \item инфимум $\inf B$ --- наибольшую нижнюю границу;
 \item супремум $\sup B$ --- наименьшую верхнюю границу.
\end{itemize}
Наименьший элемент, если существует, единственен; минимальных элементов может
быть несколько.

Конечный частичный порядок представляют диаграммой Хассе: петли и рёбра,
следующие из транзитивности, не рисуют. Линейное расширение конечного частичного
порядка всегда существует; алгоритмически его можно получить топологической
сортировкой ориентированного графа отношения покрытия.

\subsection{Функции как отношения}

Функция $f\colon A\to B$ --- специальное отношение, в котором каждому
$a\in A$ соответствует ровно один $b\in B$. Образ множества $X\subseteq A$:
\[
    f(X)=\{f(x):x\in X\},
\]
прообраз $Y\subseteq B$:
\[
    f^{-1}(Y)=\{x\in A:f(x)\in Y\}.
\]
Для прообраза всегда выполняются точные равенства
\[
 f^{-1}(Y_1\cup Y_2)=f^{-1}(Y_1)\cup f^{-1}(Y_2),
\]
\[
 f^{-1}(Y_1\cap Y_2)=f^{-1}(Y_1)\cap f^{-1}(Y_2).
\]

Функция инъективна, если $f(a_1)=f(a_2)\Rightarrow a_1=a_2$; сюръективна, если
$f(A)=B$; биективна, если одновременно инъективна и сюръективна. Обратная
функция существует тогда и только тогда, когда $f$ биективна.

\subsection{Что важно сказать на экзамене}

Основная линия ответа: множества и операции $\to$ декартово произведение $\to$
отношения $\to$ эквивалентности и порядки $\to$ функции как специальные
отношения. Центральная теорема --- соответствие отношений эквивалентности и
разбиений. Для расширенного ответа полезны мощности, диагональный аргумент
Кантора, транзитивное замыкание и различие минимального/наименьшего элемента.

\newpage
